\section{Limits, Related Work and Open Questions}\label{sec:related}

\subsection{What hashes cannot do}\label{sec:limits}\label{sec:mmv-limit}\label{sec:delay-encryption}\label{sec:batching}

\paragraph{The delay is limited by the sender's work.}
Mahmoody, Moran and Vadhan~\cite{MMV11} proved that if building a puzzle takes $N$ hash calls, an attacker with unlimited parallelism can open it in about $N$ oracle rounds. The chained puzzle needs about $N$ rounds, so it is asymptotically optimal among random-oracle constructions: no hash-only puzzle can make the sender's work grow more slowly than the delay. The attack itself makes enormous numbers of queries, so this is a limit on what a proof can promise, not a cheap way to open puzzles.

\paragraph{Delay encryption is out of reach.}
Suppose a public setup costs $S$ hash calls, after which senders encrypt to a slot with $E$ calls each and one long computation per slot opens everything. Then running the setup and encrypting is itself a puzzle that costs $S + E$ calls, so by the result above it opens in about $S + E$ rounds: with a cheap setup, the delay is limited by the cost of one encryption. A long setup is not covered by this argument, but if each slot then has a unique, quickly checkable key, computing it is a verifiable delay function, which hashes alone cannot give~\cite{MSW20,GRY25,AAKTvC26}. No method is known that opens independent puzzles for less than one full solve each.

\paragraph{Size.}
Each chain the sender runs in parallel needs its own hidden link, so a chained puzzle that lets the sender work $g$ ways in parallel has about $gb$ bytes. Whether a hash-only puzzle can be smaller is open.

\subsection{Related work}\label{sec:prior}

Rivest, Shamir and Wagner~\cite{RSW96} introduced time-lock puzzles from repeated squaring in RSA groups, which need a trusted setup or class groups and are broken by quantum computers. Mahmoody, Moran and Vadhan~\cite{MMV11} gave the chained-hash puzzle, and Afshar et al.~\cite{ACHLM23} showed that it keeps its delay against quantum openers in the quantum random-oracle model. Proofs of sequential work~\cite{MMV13,CP18} and verifiable delay functions~\cite{BBBF18} are neighbouring notions that need different things. The chained-hash idea has been implemented several times, by Branwen~\cite{Branwen11}, Johnson~\cite{Johnson13}, Todd~\cite{Todd14} and EtherDream~\cite{EtherDream24}, without a concrete bound or an explicit attacker model. Post-quantum time-lock puzzles from lattices~\cite{AMZ24} build in time logarithmic in the delay after preprocessing, but rest on circular small-secret LWE and have not been engineered. Delay encryption from isogenies~\cite{BDF21} opens a whole slot with one computation but needs a trusted setup and is broken by quantum computers. Batch-solvable puzzles~\cite{DGM24} use pairings, and \texttt{tlock}~\cite{GMR23} encrypts to a future round of a threshold beacon. Timed commitments~\cite{BN00} are the pattern \leanseal follows in the mempool. Encrypted-mempool designs are compared in \cref{sec:comparison}.

\subsection{Open questions}\label{sec:open}

\begin{enumerate}[label=\arabic*.,leftmargin=1.5em]
  \item \textbf{The speed of chained SHA-256.} Every parameter is priced against $\ell$. A circuit-level estimate from chip designers would turn the assumed advantage of 3.7 into a defensible number.
  \item \textbf{Checking puzzles.} A proof over all $nm$ chain steps would cost far more than building the puzzle. Is there a cheap check, for example one that spot-checks the chains or the links between them? Short of that, a solver that finds a malformed puzzle could prove it with a bisection game over checkpoints of the chains, which Naraghi also assumes, without analysing it, for the pads of blind simultaneous release~\cite{Naraghi26}.
  \item \textbf{A quantum bound for the exact encoding}, extending Afshar et al.\ to truncated links, the accumulator and the commitment.
  \item \textbf{Mechanism design.} Deposit sizes that survive fee inversion~\cite{Naraghi26}, solver payment, simultaneous release of keys, and the execution of late-opened transactions, tested in a protocol such as LUCID.
\end{enumerate}
